% mage-auto-heading: what-does-a-software-factory-produce
% mage-auto-heading: walking-through-a-software-factory
% Authoritative LaTeX manuscript; Markdown is a generated compatibility projection.

% <!-- part-title: DocAble Among Other Software Factories -->

% <!-- chapter-title: The Software Factory -->

\chapter{The Software Factory}

% <!-- section-terms: tolerance -->

% <!-- point: the-printer-described-a-fabricator | The printer analogy explained fabricative technology — a machine that realizes a broad range of artifacts from a description of what should be made. | terms: printer-metaphor -->

Earlier in this book, I compared coding agents to 2D and 3D printers. The comparison was about fabricative technology: unlike a specific-purpose tool such as a stapler, a fabricator can realize a broad range of artifacts from a description of what should be made.

% <!-- point: the-printer-becomes-one-piece-of-the-factory | The printer metaphor expands rather than being replaced: once fabrication becomes a means of production, the fabricator sits inside a larger system that decides what should be made, organizes production, and judges what is acceptable. | terms: printer-metaphor -->

The printer becomes one piece of a larger factory: a system that determines what should be made, organizes production, and evaluates whether the result is acceptable.

% <!-- point: scaling-fabrication-required-engineering-the-production-system | Interchangeable parts, mechanization, and machine tools each made production more repeatable and each exposed a problem, so scaling fabrication required engineering the production system around the fabricator. -->

Factories developed to make production repeatable at scales that craft production could not support. Interchangeable parts allowed components made separately to participate in the same product. Mechanization increased the amount of work that machinery could perform. Machine tools increased precision and repeatability. Each advance also exposed a problem. Faster machinery could produce defects faster than people could inspect them; a design was of little use if the available machinery could not realize it accurately; interchangeable parts were not interchangeable if different shops produced incompatible dimensions. Scaling fabrication therefore required engineering the production system around the fabricator.

Software engineering has already made much of this move. Continuous integration and delivery organize building, testing, and deployment into repeatable processes, with automated checks governing what may proceed~\citep{humble2010continuous}. Coding agents extend automation into the production of the changes themselves. The factory must now govern not only how software is built and delivered, but also how implementation decisions are delegated, what system knowledge informs them, and what evidence justifies accepting the result.

% <!-- point: process-design-and-tolerance | Two techniques from that history carry into software factories: process design determines how production proceeds, and tolerance states which variation in the product is acceptable. | terms: tolerance -->

Two ideas from manufacturing help explain what this expanded software factory must do (\autoref{process-design-and-tolerance}):

\begin{itemize}
\tightlist
\item
  \textbf{Process design} determines how production proceeds. Work is decomposed and sequenced; responsibilities, machinery, and information are arranged; measurements and decision points occur at selected stages. The process carries knowledge about how production should occur rather than requiring each fabricator to reconstruct it. Scientific management made this separation explicit by moving planning, work methods, instructions, and standards into a deliberately designed production system~\citep{taylor1911}.
\item
  \textbf{Tolerance} states which variation in the product is acceptable. A drawing can specify a dimension and its allowable variation without prescribing every motion of the machinist or every property of the finished part. Gauges and metrology then establish whether selected properties of the realized part fall within those bounds. The factory can therefore leave substantial freedom to the fabricator while retaining control over the properties that matter~\citep{winchester2018}.
\end{itemize}

As \magevalue{sec:constraints-sensors-validators-gates} showed, measurement can establish not only whether a realization satisfies a tolerance but, for quantitative properties, how much margin remains. That evidence can inform the next engineering decision without prescribing how the fabricator must work.

\begin{figure}[htbp]
\centering
\includegraphics[width=\linewidth]{book/latex/assets/process-design-and-tolerance-7f14743861.pdf}
\caption[Process design and tolerance]{\emph{Two problems of factory engineering.} Process design determines how production proceeds around the fabricator. Tolerance identifies acceptable variation in what the factory produces; measurement establishes whether a realization remains within those bounds. The two concerns are complementary: a production system may carry engineering knowledge in the organization of production, in explicit bounds on the resulting product, or in both.}
\label{process-design-and-tolerance}
\end{figure}

% <!-- point: a-better-production-system-made-a-different-artifact-practical | Watt's engine needed cylinders more accurately bored than ordinary production could manage; Wilkinson's machinery made the design practical, connecting production capability to product tolerance. | terms: tolerance -->

Watt\textquotesingle s improved steam engine shows why production capability and product tolerance must develop together. Its design required cylinders more accurately bored than ordinary eighteenth-century production could reliably make; Wilkinson\textquotesingle s precision boring machinery made cylinders sufficiently true for the engine to perform as intended. The advance was therefore more than a faster way to make an existing artifact. A more capable production system made a different artifact practical.

\section{What Does a Software Factory Produce?}\label{heading:what-does-a-software-factory-produce}

% <!-- point: software-does-not-repeatedly-produce-instances | The analogy needs one adjustment: a physical factory commonly produces many instances of a product, while a software service may be one continuously evolving system and another copy of a program costs almost nothing. -->

The analogy requires one adjustment before we apply it to software. A physical factory commonly produces many instances of a product. Software usually does not. Another copy of a program may cost almost nothing, while a software service may consist of one continuously evolving system serving millions of users.

% <!-- point: the-output-is-the-next-working-realization | The output of a software factory is not another copy of the product but the next working realization of it: each production cycle begins with a working software system and attempts to produce a changed system that still works. -->

The output of a software factory is the next working realization of the product. Software organizations repeatedly change an existing system: adding features, repairing defects, adapting to new environments, migrating dependencies, changing architectures, and responding to new requirements. Each production cycle begins with a working software system and attempts to produce a changed system that still works.

% <!-- point: a-software-factory-is-a-production-system-for-controlled-change | A software factory organizes controlled change while preserving control over consequential properties and consequences. -->

The factory\textquotesingle s job is to realize controlled change: creating and evolving software while preserving sufficient control over the properties and consequences that matter~\citep{davis2026sehandbook}. A software factory is the production system through which that change occurs.

\begin{quote}
\textbf{Software factory.} The engineered production system through which software changes are realized from engineering intent and evaluated, admitted, deployed, and improved.
\end{quote}

% <!-- point: a-factory-contains-more-than-its-fabricators | Models, tools, and controls perform different roles within the factory, whose GEE is joined by people, agents, production activities, and organizational authority. -->

Within the factory, models describe the product: architecture descriptions, service graphs, and the properties a change must preserve. Tools work on it: linters, test harnesses, validators, and analyzers. Controls bound what may be done and accepted through permissions, admission gates, and enforced policies.

The governed engineering environment developed in Chapters \magevalue{chapter:modeling|num}--\magevalue{chapter:one-problem-many-models|num} is part of this larger factory. It comprises the reusable representations, tools, procedures, and controls through which engineering work is directed and governed. The factory also includes the people and agents doing the work, the production activities they perform, and the organizational arrangements that determine what work is undertaken and who has authority over its consequences.

% <!-- point: the-two-problems-specialize-to-changes-not-copies | A software factory preserves consequential properties while changing the system itself. | terms: tolerance -->

Unlike a factory producing repeated copies of a fixed design, a software factory must preserve consequential properties while changing the system itself.

% <!-- point: representations-of-software-carry-a-carrying-cost | Representations of changing software incur construction and maintenance costs. -->

That task carries a cost that physical manufacturing can often amortize differently. A drawing for a manufactured part may remain useful across thousands of instances. A representation of a software system describes a target that is itself changing. Requirements, architecture descriptions, interface specifications, models, and tests therefore incur a carrying cost: they must evolve with the software if they are to remain useful.

% <!-- point: organizations-judge-the-carrying-cost-differently | Organizations weigh reusable representations against the cost of skilled judgment during implementation. -->
% <!-- point: the-question-was-where-judgment-is-carried-most-economically | The economic question is where consequential knowledge can be maintained and applied most economically without losing control. -->

Software organizations have made different choices about where to retain engineering judgment. Explicit representations can preserve knowledge for reuse, but require skilled labor to construct and maintain. Keeping knowledgeable engineers close to implementation avoids some of that cost, but requires them to resolve questions about behavior, structure, and tradeoffs as the software changes~\citep{kautz2018mbse}. Neither choice is intrinsically irrational. The economic question is where consequential engineering knowledge can be maintained and applied most economically without losing control.

% <!-- box-family: inset -->
% <!-- inset-domain: SOFTWARE ENGINEERING -->
% mage-box-prefix: Inset —

\begin{bookbox}{inset}{Engineering judgment and the Agile movement}
\label{inset-engineering-judgment-agile}
% <!-- point: earlier-separation-attempts-met-the-same-tradeoff | Moving decisions upstream can reduce implementation judgment but requires skilled work to construct and maintain the representations. -->
Earlier attempts to separate engineering judgment from implementation often sought to move decisions upstream. Detailed requirements, architectures, interfaces, and models would guide programmers who did not need to reconstruct the underlying design decisions. That arrangement could reduce the judgment required during implementation, but constructing and maintaining the representations demanded substantial skilled engineering labor.

% <!-- point: agile-is-intelligible-when-the-product-is-change | Agile practices placed greater weight on knowledgeable engineers close to implementation; representation costs help explain the tradeoff. -->
The Agile movement placed greater weight on keeping knowledgeable engineers close to implementation~\citep{agile-manifesto2001,scrumguide2020}. Working software became the primary measure of progress; changing requirements were expected; collaboration was preferred to elaborate handoffs; and documentation was valuable insofar as it continued to serve the work. These practices reflected an economic tradeoff. When software changes frequently, a representation that is costly to maintain may provide less value than having an engineer resolve the question directly.

The approaches need not be opposites. Both seek to produce acceptable software while controlling the cost of engineering judgment. They differ in how much knowledge is preserved in reusable representations and how much must be supplied by people during implementation. As that balance changes, so does the value of investing in the production system around the programmer.
\end{bookbox}

% <!-- point: commodity-intelligence-changes-the-tradeoff | Inexpensive capable agents can resolve ordinary implementation choices while engineers preserve selected consequential judgments and acceptance evidence. | terms: commodity-intelligence -->

Commodity intelligence changes this tradeoff. In an agentic software factory, inexpensive agents can perform substantial realization work and resolve many ordinary implementation decisions. Engineers can preserve selected consequential judgments in reusable representations and controls while leaving agents substantial freedom over realization, subject to the evidence required for acceptance.

\section{Walking Through a Software Factory}\label{heading:walking-through-a-software-factory}

% <!-- point: four-questions-orient-a-walkthrough | Four questions guide comparison: what the factory produces, what people do, how production is organized, and which product tolerances must be satisfied. | terms: tolerance -->

Four questions orient a walkthrough without prescribing how the factory should be organized:

\begin{enumerate}
\def\labelenumi{\arabic{enumi}.}
\tightlist
\item
  \textbf{What does the factory produce?} What changes enter it, and where do they originate?
\item
  \textbf{What do people do?} Which judgments and acts of realization do people retain or delegate?
\item
  \textbf{How is production organized?} How does work move through realization, evaluation, admission, and deployment, and how do lessons change that process?
\item
  \textbf{What are the tolerances?} Which properties must remain within bounds, how are those bounds represented, and what evidence establishes compliance?
\end{enumerate}

% <!-- point: software-tolerances-take-more-forms-than-a-dimension | Software tolerances include quantitative bounds and qualitative obligations, some requiring human judgment. | terms: tolerance -->

Software tolerances include quantitative bounds, invariants, architectural constraints, and qualitative obligations that may still require human judgment; each identifies variation the factory must control while leaving other choices open.

% <!-- point: the-four-questions-are-broader-than-mage | A longitudinal case explains development and mechanisms, while public industrial accounts expose variation in the allocation of work, authority, and control. -->

These questions can be answered at different levels of detail. A longitudinal account can show how a factory's arrangements developed, which problems prompted changes, and what the resulting machinery accomplished. Reports from established organizations can show how factories differ in their allocation of work, authority, and control, although they rarely expose the same development history. DocAble provides the first perspective; the industrial accounts provide the second.
